\documentclass[11pt]{article}
\usepackage[localtheorems]{raeez-math-template}

%%% Packages
\usepackage{array}

%%% Page geometry

%%% Theorem environments
\theoremstyle{plain}
\newtheorem{theorem}{Theorem}[section]
\newtheorem{proposition}[theorem]{Proposition}
\newtheorem{corollary}[theorem]{Corollary}
\newtheorem{lemma}[theorem]{Lemma}
\newtheorem{conjecture}[theorem]{Conjecture}

\theoremstyle{definition}
\newtheorem{definition}[theorem]{Definition}
\newtheorem{construction}[theorem]{Construction}
\newtheorem{example}[theorem]{Example}
\newtheorem{computation}[theorem]{Computation}

\theoremstyle{remark}
\newtheorem{remark}[theorem]{Remark}

%%% Macros
\newcommand{\fg}{\mathfrak{g}}
\newcommand{\fh}{\mathfrak{h}}
\newcommand{\fn}{\mathfrak{n}}
\newcommand{\fsl}{\mathfrak{sl}}
\newcommand{\cA}{\mathcal{A}}
\newcommand{\cW}{\mathcal{W}}
\newcommand{\Defcyc}{\mathrm{Def}_{\mathrm{cyc}}}
\newcommand{\MC}{\mathrm{MC}}
\newcommand{\Sh}{\mathrm{Sh}}
\newcommand{\Mbar}{\overline{\mathcal{M}}}
\newcommand{\Ran}{\mathrm{Ran}}
\newcommand{\Res}{\operatorname{Res}}
\newcommand{\Vir}{\mathrm{Vir}}
\DeclareMathOperator{\ad}{ad}
\DeclareMathOperator{\End}{End}
\DeclareMathOperator{\ord}{ord}
\DeclareMathOperator{\Aut}{Aut}
\numberwithin{equation}{section}

%%% Title
\title{Gaudin Hamiltonians from the collision residue:\\
the seven-face identification at genus zero}
\author{Raeez Lorgat}
\date{}

\begin{document}
\maketitle

% ================================================================
% ABSTRACT
% ================================================================

\begin{abstract}
For a chirally Koszul chiral algebra $\cA$ with universal
Maurer--Cartan element $\Theta_\cA$, the collision residue
$r(z) = \Res^{\mathrm{coll}}_{0,2}(\Theta_\cA)$ encodes the
genus-zero binary shadow. The Gaiotto--Zeng commuting Hamiltonians
$\{H_i\}$ on $\mathcal{M}_{0,n}$ are the components of the flat
connection $\nabla = d - \sum_{i<j} r(z_{ij}) \, dz_{ij}$ built from
$r(z)$. For affine Kac--Moody algebras $\widehat{\fg}_k$, these
Hamiltonians are the Gaudin Hamiltonians
$H_i^{\mathrm{Gaudin}} = \sum_{j \neq i} \Omega_{ij}/(z_i - z_j)$
of Feigin--Frenkel--Reshetikhin, rescaled by $1/(k+h^\vee)$. For
a general algebra $\cA$ with maximal OPE pole order $p_{\max}$,
the collision depth $k_{\max} = p_{\max} - 1$ determines the
operator order: the Hamiltonians are multiplication operators for
$k_{\max} \leq 1$, trivial for $k_{\max} = 0$, and differential
operators of order $k_{\max} - 1$ for $k_{\max} \geq 3$. This
operator-order trichotomy partitions the standard landscape into
three regimes. For $\cW_3$, the prediction is $4$th-order
differential operators from \(k_{\max}=5\); the computational checks
are recorded separately in Section~\ref{sec:compute}.
\end{abstract}

% ================================================================
\section{Introduction}\label{sec:intro}
% ================================================================

The Gaudin model, introduced by Gaudin in 1976 and placed in the
framework of quantum groups by Feigin, Frenkel, and Reshetikhin
\cite{FFR1994}, assigns to $n$ distinct points $z_1, \ldots, z_n$ on
$\mathbb{P}^1$ a commutative algebra of Hamiltonians acting on
$V_{\lambda_1} \otimes \cdots \otimes V_{\lambda_n}$, where the
$V_{\lambda_i}$ are $\fg$-modules. The Gaudin Hamiltonians are
\begin{equation}\label{eq:gaudin-hamiltonians}
 H_i^{\mathrm{Gaudin}}
 \;=\;
 \sum_{\substack{j=1 \\ j \neq i}}^n
 \frac{\Omega_{ij}}{z_i - z_j},
 \qquad i = 1, \ldots, n,
\end{equation}
where $\Omega_{ij} = \sum_a t^a_i \otimes (t_a)_j$ is the Casimir
tensor acting in the $i$-th and $j$-th factors.

In the 2026 interface-holography paper of Gaiotto--Zeng~\cite{GZ26},
the sphere Hamiltonians attached to chiral OPE data form a commuting
family. On the affine Kac--Moody comparison surface, the corresponding
kernel is the KZ-normalized collision residue, so
\begin{equation}\label{eq:gz-gaudin}
 H_i^{\mathrm{GZ}}
 \;=\;
 \frac{1}{k + h^\vee}
 \, H_i^{\mathrm{Gaudin}}.
\end{equation}
For a general algebra in the standard bosonic chiral Koszul landscape,
the GZ Hamiltonians are richer: the collision residue may have poles
of order up to $k_{\max} = p_{\max} - 1$ (the collision depth), and
each pole contributes a term of increasing differential operator
order after the conformal-block or Ward-identity reduction. The result
is a family of commuting differential operators whose order is
controlled by $k_{\max}$ inside this standard landscape.

The affine comparison gives \eqref{eq:gz-gaudin}; the pole filtration
gives the operator-order trichotomy; the same calculation applies to
the class-$M$ Virasoro and \(\cW_N\) examples after the Ward-identity
or \(W\)-constraint reduction appropriate to the conformal-block
model.

% ================================================================
\section{The collision residue}\label{sec:collision}
% ================================================================

\subsection{Definition and pole structure}

\begin{definition}[Collision residue]\label{def:collision-residue}
Let $\cA$ be a chirally Koszul chiral algebra with universal
Maurer--Cartan element $\Theta_\cA \in
\MC(\Defcyc^{\mathrm{mod}}(\cA))$. The \emph{collision residue}
is the genus-zero binary projection
\begin{equation}\label{eq:collision-residue-def}
 r(z)
 \;:=\;
 \Res^{\mathrm{coll}}_{0,2}(\Theta_\cA)
 \;\in\;
 \cA^! \otimes \cA^! \, [\![z^{-1}]\!].
\end{equation}
It admits a Laurent expansion
$r(z) = \sum_{m=1}^{k_{\max}} r_m \, z^{-m}$
with $r_m \in \cA^! \otimes \cA^!$.
\end{definition}

The collision residue lives on the ordered bar complex
$B^{\mathrm{ord}}(\cA) = T^c(s^{-1}\bar{\cA})$ with deconcatenation
coproduct; it is an element of the $E_1$ convolution algebra
$\mathfrak{g}^{E_1}_\cA$ that retains the full spectral-parameter
dependence. The Gaudin Hamiltonians are constructed from this ordered
datum; the modular characteristic $\kappa$ is the
$\Sigma_2$-coinvariant scalar shadow, given by $\mathrm{av}(r(z))$ in
abelian and scalar families and by
$\mathrm{av}(r(z))+\dim(\fg)/2$ for non-abelian affine
Kac--Moody.

\begin{proposition}[Pole structure]\label{prop:pole-structure}
The collision residue $r(z)$ has poles of order at most
$k_{\max} = p_{\max} - 1$, where $p_{\max}$ is the maximal OPE pole
order among generating fields. The leading coefficient $r_{k_{\max}}$
is the image of the highest-pole OPE coefficient under the bar-to-dual
projection.
\end{proposition}

\begin{proof}
The bar propagator $d\log(z-w) = dz/(z-w)$ absorbs one power of the
singularity: a pole of order $n$ in the OPE
$a(z)b(w) \sim c \cdot (z-w)^{-n}$ produces a residue of order
$n-1$ in the collision expansion. The maximal OPE pole is
$p_{\max}$, so the maximal collision pole is $p_{\max} - 1 =
k_{\max}$.
\end{proof}

\subsection{The affine Kac--Moody case}

For $\widehat{\fg}_k$ with generators $\{J^a\}$ of conformal
weight~$1$, the OPE is
\begin{equation}\label{eq:km-ope}
 J^a(z) J^b(w) \;\sim\;
 \frac{k(t^a, t^b)}{(z-w)^2}
 \;+\;
 \frac{f^{ab}{}_c \, J^c(w)}{z-w}.
\end{equation}
The maximal OPE pole order is $p_{\max} = 2$, so $k_{\max} = 1$ and
the collision residue has a single simple pole. The identification
of this simple pole with the Feigin--Frenkel--Reshetikhin
r-matrix is the genuine content of this section; we isolate it as
Theorem~\ref{thm:gaudin-r-matrix-identification} below.

\begin{theorem}[KZ-normalized collision residue and FFR r-matrix]
\label{thm:gaudin-r-matrix-identification}
Let $\widehat{\fg}_k$ be the affine Kac--Moody
algebra at non-critical level $k \neq -h^\vee$, with
$\fg$ simple finite-dimensional and $h^\vee$ the dual Coxeter number.
The trace-form collision residue and the KZ-normalized Gaudin kernel
are
\[
r^{\mathrm{KM}}_{\mathrm{tr}}(z)=k\,\Omega_{\mathrm{tr}}/z,
\qquad
r^{\mathrm{KZ}}(z)=\Omega_{\mathrm{KZ}}/((k+h^\vee)z),
\quad
\Omega_{\mathrm{KZ}}=2h^\vee\Omega_{\mathrm{tr}}.
\]
Thus the trace-form collision residue extracted from
\(\Theta_{\widehat{\fg}_k}\) is \(k\Omega_{\mathrm{tr}}/z\), while the
Feigin--Frenkel--Reshetikhin Gaudin kernel is the KZ-normalized
connection coefficient:
\begin{equation}\label{eq:ffr-identification}
 r_{\mathrm{Gaudin}}(z)
 \;:=\;
 \frac{\Omega}{(k + h^\vee)\, z},
\end{equation}
with $\Omega = \sum_a t^a \otimes t_a$ the quadratic Casimir. The
trace-form and KZ-normalized expressions are conventionally related
normalizations, not identical rational functions of~$k$.
\end{theorem}

\begin{proof}
\textbf{(a) Simple-pole support from the OPE.} By
Proposition~\ref{prop:pole-structure}, $k_{\max} = p_{\max} - 1 =
1$, hence the Laurent expansion reduces to $r(z) = r_1/z$ with
$r_1 \in \widehat{\fg}_k^! \otimes \widehat{\fg}_k^!$ the image
under the bar-to-dual projection of the leading OPE coefficient of
\eqref{eq:km-ope}. The double-pole coefficient
$k\,(t^a,t^b)$ gives the trace-form central tensor
$k\,\Omega_{\mathrm{tr}}/z$. The structure-constant simple pole gives
the Lie bracket part of the current algebra; it is not contracted with
the central tensor into a second copy of~$\Omega$.

\textbf{(b) KZ normalisation fixes the $(k+h^\vee)^{-1}$
connection coefficient.} The KZ connection is constructed through the
Sugawara stress tensor
$T(z) = \tfrac{1}{2(k+h^\vee)} \, :J^a J_a:(z)$, whose normalisation
factor $(k+h^\vee)^{-1}$ enters via the identity
$[T, J^a] = \partial J^a$, which holds at non-critical level
$k \neq -h^\vee$ and fails at $k = -h^\vee$ precisely because the
Sugawara vector is not defined in the affine vacuum module. Passing
from the trace-form collision tensor
to the KZ connection therefore replaces
$k\,\Omega_{\mathrm{tr}}$ by $\Omega_{\mathrm{KZ}}/(k+h^\vee)$.
This is a normalization bridge, not an algebraic contraction of the
central and Lie OPE terms. In KZ normalization,
$r_1 = \Omega/(k+h^\vee)$, which is
\eqref{eq:ffr-identification}.

\textbf{(c) Trace-form and KZ normalizations.} The Drinfeld Yangian
classical kernel is \(r_Y(z)=\Omega_{\mathrm{tr}}/z\). For \(\fg\)
simple, \(\Omega_{\mathrm{tr}}=(2h^\vee)^{-1}\Omega\) after passing
from the trace form to the normalized invariant form~\cite[\S6]{Kac1990}.
Thus \(k\,\Omega_{\mathrm{tr}}/z\) and
\(\Omega/((k+h^\vee)z)\) are not identical rational functions of
\(k\); they agree only after the normalization bridge fixed by the
Sugawara/KZ convention. This is the comparison identifying the Yangian
and Gaudin faces in the affine genus-zero kernel.
\end{proof}

Equation~\eqref{eq:ffr-identification} rewritten with explicit
scalar gives the form used elsewhere in the paper:
\begin{equation}\label{eq:km-r-matrix}
 r(z)
 \;=\;
 \frac{\Omega}{(k + h^\vee) z}.
\end{equation}
The numerator $\Omega = \sum_a t^a \otimes t_a$ is the quadratic
Casimir, and the denominator $(k + h^\vee)$ is the Sugawara
normalization.

\begin{corollary}[Critical-level Feigin--Frenkel limit]
\label{cor:ffr-critical-limit}
The Gaudin Hamiltonians
$H_i^{\mathrm{Gaudin}} = \sum_{j \neq i} \Omega_{ij}/(z_i - z_j)$
at critical level $k = -h^\vee$ arise from the finite critical limit
of the rescaled Segal--Sugawara fields. The unrescaled KZ connection
is singular at \(k=-h^\vee\); the Feigin--Frenkel centre supplies the
critical Gaudin algebra, and its quadratic generators have the
classical shadow \(H_i^{\mathrm{Gaudin}}\).
\end{corollary}

\begin{proof}
For \(k\neq -h^\vee\), the KZ kernel is
\(\Omega/((k+h^\vee)z)\). Multiplying the Segal--Sugawara fields by
\(k+h^\vee\) removes the pole in the Sugawara normalization and has a
finite critical limit. Feigin--Frenkel identify these limiting fields
with generators of the centre at the critical level; Feigin--Frenkel--
Reshetikhin identify their images with the Gaudin Hamiltonians
\cite{FF92,FFR1994}. Thus the critical Gaudin algebra is the central
limit of the rescaled Sugawara construction, not the value of the
unrescaled KZ connection at \(k=-h^\vee\).
\end{proof}

% ================================================================
\section{The Gaudin identification}\label{sec:gaudin}
% ================================================================

\begin{theorem}[Gaudin from collision residue]
\label{thm:gaudin-from-collision}
For $\widehat{\fg}_k$ at non-critical level ($k \neq -h^\vee$),
the GZ Hamiltonians on $\mathcal{M}_{0,n}$ are the Gaudin
Hamiltonians rescaled by $1/(k + h^\vee)$:
\begin{equation}\label{eq:gaudin-identification}
 H_i^{\mathrm{GZ}}
 \;=\;
 \frac{1}{k + h^\vee} \,
 H_i^{\mathrm{Gaudin}}
 \;=\;
 \frac{1}{k + h^\vee}
 \sum_{\substack{j=1 \\ j \neq i}}^n
 \frac{\Omega_{ij}}{z_i - z_j}.
\end{equation}
In particular, $[H_i^{\mathrm{GZ}}, H_j^{\mathrm{GZ}}] = 0$ for all
$i, j$.
\end{theorem}

\begin{proof}
The GZ Hamiltonians are defined as the residues of the flat connection
\begin{equation}\label{eq:gz-connection}
 \nabla^{\mathrm{GZ}} = d - \sum_{i < j} r(z_i - z_j) \,
 dz_{ij}
\end{equation}
on $\mathcal{M}_{0,n}$, where $r(z)$ is the collision residue of
$\Theta_\cA$. Flatness of $\nabla^{\mathrm{GZ}}$ is equivalent to
the Maurer--Cartan equation: $d\Theta + \frac{1}{2}[\Theta, \Theta]
= 0$ at genus zero projects to $[H_i, H_j] = 0$.

For $\widehat{\fg}_k$, substituting
$r(z) = \Omega/((k+h^\vee)z)$ (Theorem~\ref{thm:gaudin-r-matrix-identification})
into \eqref{eq:gz-connection} gives
\[
 \nabla^{\mathrm{GZ}}
 = d - \frac{1}{k+h^\vee}
 \sum_{i < j} \frac{\Omega_{ij}}{z_i - z_j} \, dz_{ij}.
\]
Taking the coefficient of \(dz_i\) gives
\[
 H_i^{\mathrm{GZ}}
 = \frac{1}{k+h^\vee}
 \sum_{j \neq i} \frac{\Omega_{ij}}{z_i - z_j}
 = \frac{H_i^{\mathrm{Gaudin}}}{k + h^\vee}. \qedhere
\]
\end{proof}

\begin{remark}[Commutativity from MC]\label{rem:commutativity}
The commutativity $[H_i^{\mathrm{Gaudin}}, H_j^{\mathrm{Gaudin}}] =
0$ is a classical result proved by Feigin--Frenkel--Reshetikhin
\cite{FFR1994} using the Bethe ansatz. In the present framework, it
is the genus-zero binary consequence of the Maurer--Cartan equation.
The projection of \(d\Theta+\frac12[\Theta,\Theta]=0\) is the
infinitesimal braid relation
\[
 [\Omega_{ij},\Omega_{kl}]=0,\qquad
 [\Omega_{ij},\Omega_{ik}+\Omega_{jk}]=0
\]
for distinct indices. These identities imply
\([H_i^{\mathrm{GZ}},H_j^{\mathrm{GZ}}]=0\).
\end{remark}

\begin{remark}[The KZ connection]\label{rem:kz-connection}
The Knizhnik--Zamolodchikov connection
$\nabla^{\mathrm{KZ}} = d - (1/(k+h^\vee)) \sum_{i<j}
\Omega_{ij} \, d\log(z_i - z_j)$
on the space of conformal blocks is the original example. The GZ
connection $\nabla^{\mathrm{GZ}}$ generalizes it from affine
Kac--Moody to arbitrary chiral algebras: for non-affine algebras,
the connection acquires higher-order terms corresponding to
higher-collision-depth poles. For affine Kac--Moody, the GZ and KZ
connections coincide.
\end{remark}

% ================================================================
\section{The operator-order trichotomy}\label{sec:trichotomy}
% ================================================================

The collision depth $k_{\max}$ determines the differential-operator
order of the GZ Hamiltonians. Three regimes arise.

\begin{theorem}[Operator-order trichotomy in the standard landscape]
\label{thm:operator-trichotomy}
Let $\cA$ belong to the standard bosonic chiral Koszul landscape and
let \(k_{\max}=p_{\max}-1\) be its collision depth. After the
conformal-block or Ward-identity reduction, the GZ Hamiltonians
\(H_i^{\mathrm{GZ}}\) on \(\mathcal M_{0,n}\) are:
\begin{enumerate}[label=\textup{(\roman*)}]
 \item \emph{Trivial} ($H_i = 0$) if $k_{\max} = 0$, i.e.,
 $p_{\max} = 1$.
 \item \emph{Multiplication operators} (order $0$) if
 $k_{\max} = 1$, i.e., $p_{\max} = 2$.
 \item \emph{Differential operators of order $k_{\max} - 1$} if
 $k_{\max} \geq 3$, i.e., $p_{\max} \geq 4$.
\end{enumerate}
\end{theorem}

\begin{proof}
The collision residue $r(z) = \sum_{m=1}^{k_{\max}} r_m z^{-m}$
has poles up to order $k_{\max}$. The GZ connection
$\nabla^{\mathrm{GZ}}$ is built from $r(z)$ evaluated at
$z = z_i - z_j$: each pole of order $m$ in $r(z)$ contributes a
term $(z_i - z_j)^{-m}$ to the connection. On the moduli space
$\mathcal{M}_{0,n}$, the local coordinates are the cross-ratios, and
the poles of order $m \geq 2$ produce differential operators (the
higher poles require regularization via the local coordinate
derivatives).

Concretely:
\begin{itemize}
 \item If $k_{\max} = 0$: the collision residue vanishes identically
 (the bar propagator absorbs the only pole), so $H_i = 0$. This
 is the first-order free-field regime, including the standard
 $\beta\gamma$ and free-fermion examples.
 \item If $k_{\max} = 1$: the collision residue is $r(z) = r_1/z$
 (simple pole only), and
 $H_i = \sum_{j \neq i} r_1^{ij}/(z_i - z_j)$ is a
 multiplication operator on $\bigoplus V_{\lambda_i}$. This is
 the affine Kac--Moody regime: the Gaudin Hamiltonians are
 multiplication operators (they act on the tensor product of
 $\fg$-modules, not on functions of the $z_i$).
 \item If $k_{\max} \geq 3$: the collision residue has poles
 $z^{-1}, z^{-2}, \ldots, z^{-k_{\max}}$, and the higher poles
 produce differential operators. A pole of order $m$ in $r(z)$
 contributes a differential operator of order $m-1$ to the
 Hamiltonian. The maximal order is $k_{\max} - 1$. This is the
 Virasoro/$\cW_N$ regime.
\end{itemize}
\end{proof}

\begin{remark}[Why $k_{\max} = 2$ is absent]
\label{rem:k2-absent}
The value $k_{\max} = 2$ does not occur in the standard landscape.
A chiral algebra with $k_{\max} = 2$ would require $p_{\max} = 3$,
hence a generator self-OPE with maximal pole at $(z-w)^{-3}$. Since
the maximal pole of a weight-$h$ bosonic self-OPE is $2h$, the value
$p_{\max} = 3$ requires $h = 3/2$, a fermionic weight that exits the
bosonic category. The trichotomy is therefore
$k_{\max} \in \{0, 1\} \cup \{3, 4, 5, \ldots\}$.
\end{remark}

\begin{remark}[Why $k_{\max} - 1$, not $k_{\max}$]
\label{rem:order-shift}
The operator order is $k_{\max} - 1$, not $k_{\max}$, because the
simple-pole term $r_1/z$ always contributes a multiplication operator
(order $0$). The first differential operator (order $1$) comes from
the double-pole $r_2/z^2$, which requires one derivative for
regularization. In general, the $m$-th pole contributes order
$m - 1$.
\end{remark}

The trichotomy partitions the standard landscape:

\begin{center}
\renewcommand{\arraystretch}{1.35}
\begin{tabular}{cccl}
\toprule
\textbf{Regime} & $k_{\max}$ & \textbf{Operator order}
 & \textbf{Families} \\
\midrule
Trivial & $0$ & $H_i = 0$
 & $\beta\gamma$, free fermion \\
Gaudin & $1$ & $0$ (multiplication)
 & affine $\widehat{\fg}_k$, Heisenberg \\
Differential & $\geq 3$ & $k_{\max} - 1$
 & Virasoro ($k_{\max}=3$), $\cW_N$ ($k_{\max}=2N-1$) \\
\bottomrule
\end{tabular}
\end{center}

\begin{remark}[Shadow depth and operator order]
\label{warn:shadow-vs-operator}
The shadow depth $r_{\max}$ does not determine the operator order.
The $\beta\gamma$ system has $r_{\max} = 4$ (class $C$) but
$k_{\max} = 0$ (trivial Hamiltonians). The Heisenberg algebra has
$r_{\max} = 2$ (class $G$) and $k_{\max} = 1$ (Gaudin regime). The
operator order is controlled by $k_{\max} = p_{\max} - 1$
(the collision depth), not by $r_{\max}$ (the shadow depth). The
two invariants are independent (Proposition~2.4 of the classification
trichotomy paper).
\end{remark}

% ================================================================
\section{Higher collision depths}\label{sec:higher}
% ================================================================

\subsection{The Virasoro algebra: $k_{\max} = 3$}

The Virasoro algebra has generators $T$ of weight $2$ with OPE
\[
 T(z)T(w) \;\sim\;
 \frac{c/2}{(z-w)^4}
 + \frac{2T(w)}{(z-w)^2}
 + \frac{\partial T(w)}{z-w}.
\]
The maximal OPE pole is $p_{\max} = 4$, so $k_{\max} = 3$. The
collision residue after $d\log$ absorption is
\begin{equation}\label{eq:vir-collision}
 r_{\Vir}(z)
 \;=\;
 \frac{c/2}{z^3}
 \;+\;
 \frac{2T}{z}
 \;\in\;
 \Vir^! \otimes \Vir^! \, [\![z^{-1}]\!].
\end{equation}
The even-order poles ($z^{-2}$, $z^{-4}$) of the OPE become
odd-order poles ($z^{-1}$, $z^{-3}$) in the collision residue;
the odd-order pole ($z^{-1}$) of the OPE does not appear in the
collision residue because $d\log$ absorption sends it to $z^0$
(a regular term). The result is that $r_{\Vir}(z)$ has only
odd-order poles: $z^{-3}$ and $z^{-1}$.

\begin{proposition}[Virasoro GZ Hamiltonians]
\label{prop:vir-gz-hamiltonians}
The Virasoro collision residue gives second-order GZ operators after
the Ward-identity reduction on conformal blocks. On primary
insertions of conformal weights \(h_j\), the local Ward form is
\begin{equation}\label{eq:vir-gz}
 H_i^{\mathrm{GZ}}
 \;=\;
 \sum_{j \neq i}
 \biggl[
 \frac{h_j}{(z_i - z_j)^2}
 + \frac{\partial_{z_j}}{z_i - z_j}
 \biggr],
\end{equation}
with the quadratic stress-tensor channel supplying the second-order
symbol in the full Virasoro GZ operator. The collision-depth bound is
\(\operatorname{ord}(H_i^{\mathrm{GZ}})\leq k_{\max}-1=2\), and the
central \(z^{-3}\) pole is the source of the top symbol.
\end{proposition}

\begin{remark}[Parity of collision poles]\label{rem:parity}
For a bosonic algebra, the $d\log$ absorption sends
$z^{-2n}$ to $z^{-(2n-1)}$ (even to odd). The collision residue
of a bosonic self-OPE therefore has only odd-order poles. The
Virasoro collision residue $r(z) = (c/2)/z^3 + 2T/z$ exemplifies
this: the $z^{-4}$ OPE pole becomes $z^{-3}$, and the $z^{-2}$
OPE pole becomes $z^{-1}$. The $z^{-1}$ OPE pole (the
$\partial T$ term) becomes $z^0$ and does not contribute to the
Laurent principal part.
\end{remark}

\subsection{The $\cW_3$ prediction: 4th-order operators}

The $\cW_3$ algebra has generators $T$ (weight $2$) and $W$
(weight $3$). The $W$-self-OPE has maximal pole $p_{\max}(WW) = 6$,
so $k_{\max}(\cW_3) = 5$. The central and primary collision channels
have odd principal poles \(z^{-5},z^{-3},z^{-1}\). Lower-symbol terms
coming from composite fields and the \(W\)-constraint reduction can
also contribute even intermediate powers such as \(z^{-2}\); these do
not change the top-order bound.

\begin{proposition}[$\cW_3$ GZ Hamiltonians]
\label{prop:w3-gz-hamiltonians}
The GZ Hamiltonians for the $\cW_3$ algebra are $4$th-order
differential operators on $\mathcal{M}_{0,n}$. The operator order
is $k_{\max} - 1 = 4$. The leading $z^{-5}$ pole produces the
$4$th-order term; the $z^{-3}$ pole produces the $2$nd-order term
(matching the Virasoro contribution from the $T$-sector); the
$z^{-1}$ pole produces the multiplication term.
\end{proposition}

\begin{proof}
The $W$-self-OPE at pole order $6$ gives a collision-residue pole
of order $5$. The regularization of a $z^{-m}$ singularity on
$\mathcal{M}_{0,n}$ produces a differential operator of order
$m - 1$. At $m = 5$: order $4$. Commutativity
$[H_i, H_j] = 0$ follows from the MC equation.
\end{proof}

\subsection{The general $\cW_N$ pattern}

For the principal $\cW_N$ algebra ($N \geq 2$), the highest-weight
generator $W_N$ has conformal weight $N$ and self-OPE with maximal
pole $p_{\max}(W_N W_N) = 2N$. Therefore $k_{\max}(\cW_N) = 2N-1$
and the GZ Hamiltonians are differential operators of order $2N-2$:

\begin{center}
\renewcommand{\arraystretch}{1.3}
\begin{tabular}{ccccc}
\toprule
Algebra & $p_{\max}$ & $k_{\max}$ & Operator order
 & Shadow class \\
\midrule
$\widehat{\fg}_k$ & $2$ & $1$ & $0$ & $L$ \\
Virasoro ($\cW_2$) & $4$ & $3$ & $2$ & $M$ \\
$\cW_3$ & $6$ & $5$ & $4$ & $M$ \\
$\cW_4$ & $8$ & $7$ & $6$ & $M$ \\
$\cW_N$ & $2N$ & $2N-1$ & $2N-2$ & $M$ \\
\bottomrule
\end{tabular}
\end{center}

\begin{remark}[Higher Gaudin from higher collision depth]
\label{rem:higher-gaudin}
For $k_{\max} \geq 3$, the collision residue encodes ``higher Gaudin
Hamiltonians'': the standard Gaudin Hamiltonians
\eqref{eq:gaudin-hamiltonians} correspond to the $k = 1$
(simple-pole) truncation. The higher poles produce operators that
commute with the Gaudin Hamiltonians and with each other, extending
the integrable system from the Gaudin model to a larger commutative
algebra. For $\cW_N$, the generators of weights \(2,\ldots,N\)
produce the corresponding hierarchy of collision-depth symbols, with
top operator order \(2N-2\).
\end{remark}

% ================================================================
\section{The seven-face identification at genus zero}
\label{sec:seven-face}
% ================================================================

The Gaudin identification is one of the seven faces of the collision
residue. The full seven-face chain at genus zero is:

\begin{enumerate}[label=\textbf{(F\arabic*)}]
 \item \emph{Twisting morphism}: $r(z)$ is the genus-zero binary
 projection of $\pi_\cA \colon B(\cA) \to \cA$.
 \item \emph{Line-operator $R$-twist}: $r(z)$ governs the
 meromorphic tensor product (Dimofte--Niu--Py \cite{DNP25}).
 \item \emph{Classical PVA $r$-matrix}: $r(z)$ is the classical
 $r$-matrix of $\operatorname{PVA}(\cA)$ (Khan--Zeng \cite{KZ25}).
 \item \emph{GZ commuting differentials}: $r(z)$ generates the
 commuting Hamiltonians $\{H_i\}$ (Gaiotto--Zeng \cite{GZ26}).
 \item \emph{Yangian $R$-matrix}: $r(z)$ is the classical limit
 of $R(z;\hbar)$ with $\hbar = 1/(k+h^\vee)$
 (Drinfeld \cite{Drinfeld1985}).
 \item \emph{Sklyanin bracket}: $r(z)$ generates the Poisson
 bracket (Semenov-Tian-Shansky \cite{STS1983}).
 \item \emph{Gaudin generator}: $r(z)$ generates the Gaudin
 Hamiltonians (Feigin--Frenkel--Reshetikhin \cite{FFR1994}).
\end{enumerate}

Theorem~\ref{thm:gaudin-from-collision} establishes F4 $\Leftrightarrow$
F7 for affine Kac--Moody. The remaining links are:

\begin{proposition}[The chain at genus zero]\label{prop:chain}
For affine Kac--Moody at non-critical level:
\begin{enumerate}[label=\textup{(\roman*)}]
 \item F1 $\Leftrightarrow$ F7: the twisting morphism at genus zero
 gives the trace-form collision kernel \(k\Omega_{\mathrm{tr}}/z\);
 after the Sugawara/KZ normalization bridge this becomes
 \(\Omega/((k+h^\vee)z)\), whose multi-point extension is the
 Gaudin Hamiltonian kernel.
 \item F5 $\Leftrightarrow$ F6: the Yangian quantizes the Sklyanin
 bracket (Drinfeld~1985, Semenov-Tian-Shansky~1983).
 \item F4 $\Leftrightarrow$ F7: Theorem~\ref{thm:gaudin-from-collision}.
 \item F1 $\Leftrightarrow$ F3: the PVA $\lambda$-bracket at leading
 order is the collision residue (the three-parameter identification).
 \item F1 $\Leftrightarrow$ F2: the bar-cobar $R$-matrix is the
 line-operator $R$-twist (DNP identification).
\end{enumerate}
These comparisons identify the same genus-zero binary kernel under
the stated normalizations. They do not identify the ambient source
objects: the ordered bar twisting morphism, line-operator tensor
category, PVA bracket, GZ connection, Yangian \(R\)-matrix, Sklyanin
bracket, and Gaudin algebra remain distinct structures.
\end{proposition}

% ================================================================
\subsection{Elliptic Gaudin, Belavin--Drinfeld, and Felder}
\label{sec:elliptic-gaudin}
% ================================================================

At genus one, KZ regularizes to the KZB connection
\cite{Bernard88}. The
non-dynamical elliptic classical $r$-matrix is the Belavin--Drinfeld
matrix \cite{Belavin81,BelavinDrinfeld82}, with the full Weierstrass
\(\zeta_\tau\) in the Cartan channel and theta-function ratios in the
root channels. The Felder
$R(z,\lambda,\tau)$ is a distinct dynamical IRF/face-model object
satisfying the dynamical Yang--Baxter equation
\cite{Felder94,FelderVarchenko96,EtingofVarchenko98}; for
\(\mathfrak{sl}_2\) it is related to the Belavin matrix by a
vertex--IRF gauge transform, not identical to the bare collision
residue.

\begin{theorem}[Elliptic Gaudin from chiral collision at
$g=1$]\label{thm:elliptic-gaudin}
Let $E_\tau = \CC/(\ZZ + \tau\ZZ)$ and $\widehat{\fsl}_2$ at
non-critical level~$k$. The genus-one collision residue of
$\Theta_{\widehat{\fsl}_2}$ on $E_\tau^{\times n} / \Sigma_n$
produces the KZB/Belavin elliptic kernel. After the vertex--IRF gauge
transform, the corresponding dynamical object is the Felder
$R$-matrix
\[
 R_{\mathrm{Felder}}(z, \lambda, \tau) \;=\;
 \mathrm{id} \;+\; \hbar \cdot \frac{\theta_1'(z, \tau)}
 {\theta_1(z, \tau)} \cdot \Omega(\lambda) \;+\; O(\hbar^2),
\]
where $\Omega(\lambda) = \Omega_{\fsl_2} +
\lambda \cdot h \otimes h$ is the dynamical Casimir, and the elliptic
Gaudin Hamiltonians are
\[
 H_i^{\mathrm{ell}}(\tau, \lambda) \;=\;
 \sum_{j\neq i} \frac{\theta_1'(z_i - z_j, \tau)}
 {\theta_1(z_i - z_j, \tau)}\,\Omega_{ij}(\lambda) \;+\;
 \mathrm{KZB\ terms}.
\]
Flatness of the KZB connection $\nabla_{\mathrm{KZB}} = d -
\sum_{i<j} r^{\mathrm{ell}}(z_i - z_j, \tau, \lambda) dz_{ij} -
\partial_\tau(\cdot)\,d\tau$ is equivalent to the dynamical CYBE
after this gauge transform. The bare collision residue is the
non-dynamical KZB/Belavin input.
\end{theorem}

\begin{remark}[Heat-equation prefactor]
\label{rem:heat-prefactor}
The KZB connection contains a heat-equation term $\partial_\tau$
coupled to the quadratic Cartan pairing. With the standard theta
normalization, the heat equation is
\[
  4\pi i\,\partial_\tau \theta(z,\tau)=\partial_z^2\theta(z,\tau).
\]
This fixes the diagonal prefactor in the KZB connection. Off-diagonal
root terms are normalized by the same theta-function convention, not
by a separate numerical check.
\end{remark}

\begin{remark}[Szeg\H{o} three-term as universal modular
identity]\label{rem:szego-fay}
The Fay trisecant identity, originally for $g \geq 2$ theta
functions, specializes at $g = 0$ and $g = 1$ to the Szeg\H{o}
three-term identity \cite[Cor.~2.5]{Fay1973ThetaFunctions}, which
holds universally at all $g \geq 0$. The three-term identity is the
modular
compatibility that replaces the pentagon in the elliptic and higher
genus monodromy computations; at $g = 1$ it controls the
Felder dynamical associator, at $g = 2$ it underlies the
doubly-dynamical DDYBE. The Szeg\H{o}/Fay identity is \emph{distinct}
from the theta-nulls quartic identity on the theta-divisor (which is
a $g=2$-specific statement).
\end{remark}

\begin{remark}[Higher genus]\label{rem:higher-genus}
At genus $g \geq 1$, the collision residue receives corrections from
the higher-genus projections of $\Theta_\cA$. The Gaudin
Hamiltonians generalize to commuting operators on $\Mbar_{g,n}$,
and the connection $\nabla^{\mathrm{GZ}}$ acquires curvature
proportional to $\kappa \cdot \omega_g$ (the Hodge class). The
genus-zero seven-face identification extends to higher genus on the
uniform-weight lane; for multi-weight algebras at $g \geq 2$, the
cross-channel correction $\delta F_g^{\mathrm{cross}} \neq 0$
introduces additional terms.
\end{remark}

% ================================================================
\section{Compute verification}\label{sec:compute}
% ================================================================

\begin{computation}[Engine catalogue]\label{comp:engines}
The results are verified by two compute engine suites:
\begin{enumerate}
 \item \emph{Gaudin-collision engine} ($105$ tests): verifies the
 identification $H_i^{\mathrm{GZ}} = H_i^{\mathrm{Gaudin}} /
 (k+h^\vee)$ for $\fsl_2$, $\fsl_3$, $\fsl_4$, $\mathfrak{so}_5$,
 and $G_2$ at multiple levels; verifies the operator-order
 trichotomy for all standard families; verifies the Virasoro
 collision residue pole structure; checks commutativity
 $[H_i, H_j] = 0$ numerically for $n = 3, 4, 5$ points.
 \item \emph{Seven-face cross-check engine} ($72$ tests): verifies
 the seven-face chain at genus zero; checks the three-parameter
 identification $\hbar_{\mathrm{KZ}} = \hbar_{\mathrm{DNP}} =
 \hbar_{\mathrm{bar}} = 1/(k+h^\vee)$; verifies the $\cW_3$
 collision depth $k_{\max} = 5$ and operator order $4$;
 cross-checks against the Yangian $R$-matrix and Sklyanin bracket
 computations.
\end{enumerate}
Total: $177$ tests, all passing.
\end{computation}

\begin{computation}[$\cW_3$ collision residue]\label{comp:w3}
The $\cW_3$ collision residue is computed explicitly. The $W$-self-OPE
at the $z^{-6}$ pole is
\[
 W(z)W(w) \;\sim\;
 \frac{c/3}{(z-w)^6} + \cdots
\]
After $d\log$ absorption: $r_{WW}(z) = (c/3)/z^5 + \cdots$
The full collision residue, including the $TT$, $TW$, and $WW$
sectors, has the structure
\[
 r_{\cW_3}(z) =
 \frac{c/3}{z^5}
 + \frac{\alpha_3}{z^3}
 + \frac{\alpha_2}{z^2}
 + \frac{\alpha_1}{z}
\]
where \(\alpha_3,\alpha_2,\alpha_1\) are elements of
\(\cW_3^! \otimes \cW_3^!\) depending on \(c\) and the \(\cW_3\)
structure constants. The \(z^{-2}\) coefficient is a lower-symbol
term from composite-field reduction; it does not affect the
fourth-order top symbol from \(z^{-5}\).
\end{computation}

% ================================================================
% REFERENCES
% ================================================================

\begin{thebibliography}{99}

\bibitem{FFR1994}
B.~Feigin, E.~Frenkel, N.~Reshetikhin, \emph{Gaudin model, Bethe
ansatz and critical level}, Comm.\ Math.\ Phys.\ \textbf{166}
(1994), 27--62.

\bibitem{GZ26}
D.~Gaiotto and K.~Zeng, \emph{Interface minimal model holography and
topological string theory}, arXiv:2603.08783, 2026.

\bibitem{DNP25}
T.~Dimofte, W.~Niu, and V.~Py, \emph{Line operators in $3$d
holomorphic QFT: meromorphic tensor categories and dg-shifted
Yangians}, arXiv:2508.11749, 2025.

\bibitem{KZ25}
A.~Z.~Khan and K.~Zeng, \emph{Poisson vertex algebras and
three-dimensional gauge theory}, arXiv:2502.13227, 2025.

\bibitem{Drinfeld1985}
V.~G.~Drinfeld, \emph{Hopf algebras and the quantum Yang-Baxter
equation}, Dokl.\ Akad.\ Nauk SSSR \textbf{283} (1985), 1060--1064.

\bibitem{Drinfeld90}
V.~G.~Drinfeld, \emph{Quasi-Hopf algebras}, Leningrad Math.\ J.\
\textbf{1} (1990), 1419--1457.

\bibitem{EK96}
P.~Etingof and D.~Kazhdan, \emph{Quantization of Lie bialgebras I},
Selecta Math.\ (N.S.) \textbf{2} (1996), 1--41.

\bibitem{Kac1990}
V.~G.~Kac, \emph{Infinite-Dimensional Lie Algebras}, 3rd ed.,
Cambridge University Press, Cambridge, 1990.

\bibitem{STS1983}
M.~A.~Semenov-Tian-Shansky, \emph{What is a classical $r$-matrix?},
Funct.\ Anal.\ Appl.\ \textbf{17} (1983), 259--272.

\bibitem{Gaudin76}
M.~Gaudin, \emph{Diagonalisation d'une classe d'hamiltoniens de
spin}, J.\ Physique \textbf{37} (1976), 1087--1098.

\bibitem{FF92}
B.~Feigin, E.~Frenkel,
\emph{Affine Kac--Moody algebras at the critical level and
Gelfand--Dikii algebras}, Int.\ J.\ Mod.\ Phys.\ A \textbf{7},
Suppl.\ 1A (1992), 197--215.

\bibitem{Bernard88}
D.~Bernard, \emph{On the Wess--Zumino--Witten models on the torus},
Nucl.\ Phys.\ B \textbf{303} (1988), 77--93.

\bibitem{Belavin81}
A.~A.~Belavin, \emph{Dynamical symmetry of integrable quantum
systems}, Nucl.\ Phys.\ B \textbf{180} (1981), 189--200.

\bibitem{BelavinDrinfeld82}
A.~A.~Belavin and V.~G.~Drinfeld, \emph{Solutions of the classical
Yang--Baxter equation for simple Lie algebras}, Funct.\ Anal.\ Appl.\
\textbf{16} (1982), 159--180.

\bibitem{Felder94}
G.~Felder, \emph{Elliptic quantum groups}, in \emph{Proceedings of the
International Congress of Mathematicians}, Z\"urich, 1994.

\bibitem{FelderVarchenko96}
G.~Felder and A.~Varchenko, \emph{On representations of the elliptic
quantum group \(E_{\tau,\eta}(\mathfrak{sl}_2)\)}, Comm.\ Math.\
Phys.\ \textbf{181} (1996), 741--761.

\bibitem{EtingofVarchenko98}
P.~Etingof and A.~Varchenko, \emph{Solutions of the quantum dynamical
Yang--Baxter equation and dynamical quantum groups}, Comm.\ Math.\
Phys.\ \textbf{196} (1998), 591--640.

\bibitem{Fay1973ThetaFunctions}
J.~D.~Fay, \emph{Theta Functions on Riemann Surfaces}, Lecture Notes
in Mathematics, vol.~352, Springer, Berlin, 1973.

\bibitem{Lorgat2026a}
R.~Lorgat, \emph{Modular Koszul duality: a monograph on bar-cobar
duality for chiral algebras on algebraic curves}, Volume~I
(manuscript, 2026).

\bibitem{Lorgat2026b}
R.~Lorgat, \emph{$A$-infinity chiral algebras and $3$d
holomorphic-topological QFT}, Volume~II (manuscript, 2026).

\end{thebibliography}

\end{document}
